\documentclass[11pt]{article}
\usepackage[margin=2.5cm]{geometry}
\usepackage{amsmath,amssymb}
\usepackage{booktabs}
\usepackage{array}
\usepackage{xcolor}
\usepackage{hyperref}
\usepackage{listings}
\usepackage{microtype}

\newcommand{\E}{\mathbb{E}}
\newcommand{\Wx}{W\!x}
\newcommand{\code}[1]{\texttt{#1}}
\let\origunderscore\_
\renewcommand{\_}{\origunderscore\allowbreak} % let long code_identifiers wrap at underscores instead of overflowing table columns
\newcommand{\verified}{\textcolor{blue!50!black}{[verified from code]}}
\newcommand{\fixed}{\textcolor{green!40!black}{[fixed 2026-08-31]}}
\newcommand{\suspect}{\textcolor{red!60!black}{[does not match Eq.~\ref{eq:paper-norm}, unresolved]}}
\newcommand{\none}{\textcolor{gray}{[none implemented]}}
\newcommand{\notefix}{\textcolor{orange!70!black}{[note corrected 2026-08-31 -- code unchanged]}}

\lstset{basicstyle=\ttfamily\small, breaklines=true}
\emergencystretch=3em

\title{Catalog of 1D potentials in \texttt{codes/potentials/potentials\_1d.py}}
\author{Generated with Claude Code, from a code-reading session on 2026-08-31}
\date{}

\begin{document}
\maketitle

\section*{Scope and notation}

This note catalogs every potential class in \code{codes/potentials/potentials\_1d.py},
in file order, from \code{Scattering\_First\_Order\_1d} up to and including
\code{L2p1\_norm}. \code{Scalar\_GGD\_KRegion} and everything after it (the
generalized-Gaussian fitting potentials) are out of scope. Three infrastructure
classes precede the catalog (\code{Potential}, \code{Potential\_Prepare},
\code{Potential\_Parallel}) — they are base/wrapper classes with no mathematical
content of their own and are not covered here.

\paragraph{Filter banks.} \code{get\_1d\_potentials} (\code{codes/potentials\_builder.py})
builds every potential from one of three filter banks, all constructed by
\code{return\_Filters} (\code{codes/filters\_bank.py:21}):
\begin{itemize}
  \item \textbf{\code{filters}} — the \emph{coarse} bank, $Q=1$ wavelet per octave.
        \textbf{Always} includes one low-pass channel $\phi$ concatenated as its
        last channel, alongside $J$ band-pass channels $\psi_1,\dots,\psi_J$: $J+1$
        channels total.
  \item \textbf{\code{filters\_Q}} — the \emph{fine} bank, $Q$ wavelets per octave
        (whatever $Q$ the run was configured with, e.g. $Q=3$). Same structure:
        $JQ$ band-pass channels plus the same one low-pass channel concatenated
        at the end: $JQ+1$ channels total.
  \item \textbf{\code{filters\_Phi}} — a \emph{separate} stack of $J$ different
        low-pass filters, one per scale $j=1,\dots,J$ (increasingly smoothed).
        Not a duplicate of the low-pass channel bundled into \code{filters}/\code{filters\_Q}.
\end{itemize}
Several call sites slice off the bundled low-pass channel explicitly with
\code{[:,:-1]} when a potential should only see band-pass channels.

\paragraph{Wavelet coefficient.} $\Wx[\lambda](t) := (\psi_\lambda * x)(t)$, a
complex-valued coefficient at scale/channel $\lambda$ and time $t$ (Morlet
wavelets are complex). $\E[\cdot]$ denotes the empirical average over batch and
time, $\frac{1}{BT}\sum_{b,t}$, unless stated otherwise. $\sigma(\cdot)$ denotes
the logistic sigmoid; $\mathrm{abs\_eps}(z) = \sqrt{\mathrm{Re}(z)^2+\mathrm{Im}(z)^2+\varepsilon}$
(\code{utils\_potentials.py:135}) is used with $\varepsilon=0$ everywhere in this
file, i.e. it is exactly $|z|$ in every call site checked.

\paragraph{The micro-normalization mechanism (\code{fit\_micro}/\code{self.norm}).}
Several classes below define \code{fit\_micro(x)}: called \emph{once}, on real
data, before fitting or generation (never on synthetic samples). It sets
\code{self.norm}, a per-scale rescaling factor that \code{forward()} divides the
\emph{raw first-layer coefficient} by, before any further nonlinearity. The
motivation (from the scattering-spectra literature): this codebase's data has a
power-law spectrum, so $\E[|\Wx_s|^2]$ varies hugely across scales $s$; without
correcting for that, whichever scale carries the most raw power dominates the
fitted $\theta$ purely because it is louder, and fitting becomes ill-conditioned
across frequencies. The reference normalization, generalizing Eq.~22 of the
scattering-spectra paper (Morel et al.) — divide the wavelet coefficient by its
own standard deviation once, and every statistic built from it inherits the
correctly-cancelling power of $\sigma$:
\begin{equation}
  \label{eq:paper-norm}
  \Wx[\lambda] \;\longmapsto\; \Wx[\lambda]/\sigma[\lambda], \qquad
  \sigma[\lambda] := \big(\E[|\Wx[\lambda]|^2]\big)^{1/2}.
\end{equation}
Each subsection below states explicitly what \code{self.norm} actually is, and
whether it matches \eqref{eq:paper-norm} (verified, or fixed on 2026-08-31 during
this session), or departs from it (flagged, unresolved).

\bigskip
\hrule
\bigskip

\section{\code{Scattering\_First\_Order\_1d}}
\textbf{Code:} \code{potentials\_1d.py:81--106}. \quad
\textbf{Term:} \code{'Scattering\_First\_Order'}.

\paragraph{Formula.}
\begin{equation}
  S_1[\lambda] = \E\big[\,|\Wx[\lambda](t)|\,\big], \qquad \lambda = 1,\dots,JQ.
\end{equation}
The plain first absolute moment (modulus) of the wavelet coefficient, per fine scale.

\paragraph{Filter bank.} \code{potentials\_builder.py:190}: \code{Scattering\_First\_Order\_1d(filters\_Q[:,:-1])}
— the \textbf{fine} bank, \textbf{band-pass channels only} (the bundled low-pass
is sliced off). Runs at whatever $Q$ the experiment was configured with (e.g.
$Q=3$), \emph{not} $Q=1$. \verified

\paragraph{Normalization.} \none. No \code{fit\_micro} override — inherits the
base \code{Potential.fit\_micro}, a no-op. \code{self.norm} is never set, and
\code{forward()} never references it. This statistic is exactly as
\eqref{eq:paper-norm} says it should \emph{not} be left: fully exposed to
whatever power each scale $\lambda$ happens to carry.

\bigskip
\hrule
\bigskip

\section{\code{Scattering\_Second\_Order\_1d}}
\textbf{Code:} \code{potentials\_1d.py:114--137}. \quad
\textbf{Terms:} \code{'Scattering\_Second\_Order'}, \code{'Scattering\_Second\_Order\_Q1'}.

\paragraph{Formula.}
\begin{equation}
  S_2[\lambda] = \E\big[\,\Wx[\lambda](t)\,\overline{\Wx[\lambda](t)}\,\big]
              = \E\big[\,|\Wx[\lambda](t)|^2\,\big].
\end{equation}
The plain second moment (power spectrum) of the wavelet coefficient, per scale.

\paragraph{Filter bank.} Instantiated \emph{twice} in \code{potentials\_builder.py}:
\begin{itemize}
  \item Line 193: \code{Scattering\_Second\_Order\_1d(filters\_Q)} — \textbf{fine}
        bank, term \code{'Scattering\_Second\_Order'}. \emph{Includes} the bundled
        low-pass channel (no slicing here, unlike the first-order sibling above).
  \item Line 211: \code{Scattering\_Second\_Order\_1d(filters)} — \textbf{coarse}
        ($Q=1$) bank, term \code{'Scattering\_Second\_Order\_Q1'}. Also includes
        the low-pass channel.
\end{itemize}
\verified

\paragraph{Normalization.} \none. Same as \code{Scattering\_First\_Order\_1d}:
no \code{fit\_micro}, no \code{self.norm} division.

\bigskip
\hrule
\bigskip

\section{\code{Scattering\_Second\_Order\_Bulk\_1d}}
\textbf{Code:} \code{potentials\_1d.py:139--213}. \quad
\textbf{Term:} \code{'Scattering\_Second\_Order\_bulk'} (comment in code:
``experiment for lagrangian turbulence''; only on the distribution bulk).

\paragraph{Formula.} A \emph{trimmed} second moment: the same $|\Wx[\lambda](t)|^2$
as above, but weighted by a smooth window that suppresses the tail of the
distribution of $|\mathrm{Re}(\Wx[\lambda](t))|$ (note: the window is built from
the \emph{real part}'s magnitude, not the complex modulus, while the quantity
being averaged \emph{is} the full complex modulus squared). Calibrated once, from
data, in \code{fit(x)}:
\begin{align}
  c[\lambda] &= \text{the \code{bulk\_quantile} (default 0.90) quantile of } |\mathrm{Re}(\Wx[\lambda](t))| \text{ over all } (b,t),\\
  s[\lambda] &= \tfrac{1}{4}\big(Q_{0.95}(|\mathrm{Re}(\Wx[\lambda])|) - Q_{0.85}(|\mathrm{Re}(\Wx[\lambda])|)\big), \text{ clamped to } \ge 10^{-6}.
\end{align}
Then, per-sample:
\begin{equation}
  w_{\text{bulk}}[\lambda](t) = 1 - \sigma\!\left(\frac{|\mathrm{Re}(\Wx[\lambda](t))| - c[\lambda]}{s[\lambda]}\right), \qquad
  S_2^{\text{bulk}}[\lambda] = \E_t\big[\,w_{\text{bulk}}[\lambda](t)\cdot|\Wx[\lambda](t)|^2\,\big].
\end{equation}
A robust/trimmed power spectrum: samples whose real part exceeds the bulk
quantile threshold are smoothly down-weighted rather than hard-clipped.

\paragraph{Filter bank.} \code{potentials\_builder.py:197}: \code{Scattering\_Second\_Order\_Bulk\_1d(filters\_Q)}
— \textbf{fine} bank, full (band-pass + low-pass). \verified

\paragraph{Normalization.} Not a \code{fit\_micro}/\code{self.norm}-style
rescaling of the output — instead uses its own \code{fit(x)} method to calibrate
the bulk/tail cutoff $(c[\lambda], s[\lambda])$ directly from data. The output
itself is never divided by anything. This addresses a different problem
(robustness to outliers/tails) than \eqref{eq:paper-norm} (comparability across
scales' own power) — the two are not interchangeable, and this class still has
none of the latter.

\bigskip
\hrule
\bigskip

\section{\code{Scattering\_Third\_Order\_Real\_1d}}
\textbf{Code:} \code{potentials\_1d.py:215--269}. \quad
\textbf{Term:} \code{'Scattering\_Third\_Order\_Real'}.
(\code{'Scattering\_Third\_Order\_Imag'} is wired to a \code{pass} in
\code{potentials\_builder.py:203} — the corresponding \code{Scattering\_Third\_Order\_Imag\_1d}
class is commented out and does not exist as a usable term.)

\paragraph{Formula.} A genuine \emph{cross-scale} correlation — the classic
scattering third-order coefficient, correlating the coefficient \emph{at the
second-transform's own scale} $\mu$ against a second wavelet transform, built
at that same scale $\mu$, of the modulus envelope from a finer (or equal)
scale $\lambda \le \mu$:
\begin{equation}
  S_3[\lambda,\mu] = \E_t\Big[\,\mathrm{Re}\Big(
     \overline{\Wx[\mu](t)} \cdot \big(\psi_{\mu} * |\Wx[\lambda](\cdot)|\big)(t)
  \Big)\Big], \qquad \mu \ge \lambda.
\end{equation}
\notefix\ Traced by hand from \code{forward()} (\code{potentials\_1d.py:222--235})
and confirmed numerically against dummy tensors, since the broadcasting there
is easy to misread: the coefficient that gets conjugated is
$\overline{\Wx[\mu]}$ (same scale as the second-layer filter). This is the standard scattering-transform convention (a
fine-scale envelope, re-filtered at a coarser scale $\mu$, correlated against
the coefficient. 

\paragraph{Filter bank and a real caveat at $Q>1$.} \code{potentials\_builder.py:200}:
\code{Scattering\_Third\_Order\_Real\_1d(J, filters\_Q)} — \textbf{fine} bank,
full (no slicing), with \code{self.J = J} (\code{args.J}, the base scale count,
\emph{not} \code{J*Q}) and \code{self.filters = filters\_Q} ($JQ+1$ channels).
\verified\ The index pairs $(\lambda,\mu)$ are generated by
\code{indices\_third\_order(self.J, 1)}, which builds them assuming
\code{num\_filters = self.J+1} — \emph{i.e. it only ever emits indices in
$[0, J]$}, regardless of how many channels the bank passed to \code{forward()}
actually has. \notefix\ I confirmed this numerically
(\code{indices\_third\_order(J,1)} with $J{=}8,Q{=}3$ tops out at index 8, never
9--24) rather than trusting the shapes by eye.
\begin{itemize}
  \item At \textbf{$Q=1$} (\code{filters\_Q} then equals the coarse bank, $J+1$
        channels) this is exactly right: every channel of the bank passed in is
        reachable, and $\mu$ correctly ranges up to and including the low-pass
        channel. This is the regime this term is actually used in.
  \item At \textbf{$Q>1$} (the fine bank actually has $JQ+1$ channels) this is
        \textbf{broken}: channels $J{+}1,\dots,JQ$ of \code{filters\_Q} are
        silently never referenced by this potential at all. The term still
        runs without error and returns a plausible-looking tensor of the shape
        \code{self.num\_coefficients} $= J((1+J)/2+1)$ predicts — it just
        quietly ignores most of the fine bank's scales rather than computing
        the cross-scale statistic over all of them. \textbf{Do not use
        \code{'Scattering\_Third\_Order\_Real'} with $Q>1$ filters as currently
        implemented} — fixing it would mean calling
        \code{indices\_third\_order} with an argument reflecting the actual
        $JQ+1$ channel count, not \code{self.J}.
\end{itemize}

\paragraph{Normalization.} \none. No \code{fit\_micro} defined.

\bigskip
\hrule
\bigskip

\section{\code{Scattering\_Fourth\_Order\_Real\_1d}}
\label{sec:s4real}
\textbf{Code:} \code{potentials\_1d.py:273--398} (post-fix). \quad
\textbf{Terms:} \code{'Scattering\_Fourth\_Order\_Real'}, \code{'Scattering\_Fourth\_Order\_Real\_Q1'}.

\paragraph{Formula.} Correlates the second-layer (coarse-scale) wavelet
transform of the \emph{modulus} envelope $|\Wx[s]|$ at two fine scales
$s_1 \le s_2$, at a shared coarse scale $j$:
\begin{align}
  \Wx^{\text{norm}}[s](t) &:= \Wx[s](t)/\sigma[s], \qquad
  E[s](t) := |\Wx^{\text{norm}}[s](t)|, \\
  W_j E[s](t) &:= (\phi_j * E[s](\cdot))(t) \quad \text{(second, coarse-bank transform)}, \\
  S_4^{\text{Re}}[s_1,s_2,j] &= \E_t\Big[\,\mathrm{Re}\big(W_j E[s_1](t)\cdot \overline{W_j E[s_2](t)}\big)\Big].
\end{align}
$j$ ranges over \emph{all} coarse channels including the low-pass one
(\code{include\_lowpass} defaults to \code{True} for this class — not passed
explicitly at any call site). Only $s_1\le s_2$ is stored: $\mathrm{Re}(z)$ is
symmetric under swapping $s_1\!\leftrightarrow\!s_2$, so the other half is
redundant.

\paragraph{Filter bank.} Two wirings, both in \code{potentials\_builder.py}:
\begin{itemize}
  \item Line 207: \code{Scattering\_Fourth\_Order\_Real\_1d(J, Q, filters, filters\_Q[:,:-1])}
        — term \code{'Scattering\_Fourth\_Order\_Real'}: second layer on the
        \textbf{coarse} bank, first layer on the \textbf{fine} bank (band-pass
        only), at whatever $Q$ was configured.
  \item Line 216: \code{Scattering\_Fourth\_Order\_Real\_1d(J, 1, filters, filters[:,:-1])}
        — term \code{'Scattering\_Fourth\_Order\_Real\_Q1'}: both layers on the
        \textbf{coarse} ($Q=1$) bank explicitly.
\end{itemize}
\verified

\paragraph{Normalization.} \code{self.norm}$[s] = \big(\E[|\Wx[s]|^2]\big)^{0.5} = \sigma[s]$
\fixed — matches \eqref{eq:paper-norm} exactly. Previously
\code{self.norm}$[s]=\big(\E[|\Wx[s]|^2]\big)^{-1/4}$; the old line is kept
commented out in the source. Derivation of why $0.5$ is the exponent that
cancels the scale-dependence of $S_4^{\text{Re}}$ (independent of whether the
envelope is $|\Wx|$ or $|\Wx|^2$, see Section~\ref{sec:s4mod2real}): tracking how
the raw (unnormalized) statistic scales with $E_{s_1}:=\E[|\Wx[s_1]|^2]$ and
$E_{s_2}$ gives a factor $E_{s_1}^{1/2-p}E_{s_2}^{1/2-p}$ for exponent $p$ on
$\sigma[s]$ in the denominator; this vanishes (scale-invariance) only at $p=1/2$.

\bigskip
\hrule
\bigskip

\section{\code{Scattering\_Fourth\_Order\_Imag\_1d}}
\textbf{Code:} \code{potentials\_1d.py:400--513}. \quad
\textbf{Terms:} \code{'Scattering\_Fourth\_Order\_Imag'}, \code{'Scattering\_Fourth\_Order\_Imag\_Q1'}.

\paragraph{Formula.} Identical construction to Section~\ref{sec:s4real}, keeping
the imaginary part instead:
\begin{equation}
  S_4^{\text{Im}}[s_1,s_2,j] = \E_t\Big[\,\mathrm{Im}\big(W_j E[s_1](t)\cdot \overline{W_j E[s_2](t)}\big)\Big], \qquad s_1 < s_2 \text{ strictly}.
\end{equation}
$\mathrm{Im}(z)$ is \emph{antisymmetric} under $s_1\!\leftrightarrow\!s_2$
(swapping gives $\mathrm{Im}(\bar z)=-\mathrm{Im}(z)$), so the diagonal
$s_1=s_2$ is identically zero and is always excluded (\code{offset=1}, not
optional here — unlike the Real class, this class has no \code{include\_diag}
constructor argument at all). \textbf{Also}: \code{include\_lowpass=False} is
hardcoded in this class's index construction, so $j$ ranges only over the $J$
band-pass coarse channels, \emph{excluding} the low-pass one — a second,
independent difference from the Real class (Section~\ref{sec:s4real}) beyond
the diagonal exclusion.

\paragraph{Filter bank.} Two wirings:
\begin{itemize}
  \item Line 219: \code{Scattering\_Fourth\_Order\_Imag\_1d(J, Q, filters, filters\_Q[:,:-1])}
        — term \code{'Scattering\_Fourth\_Order\_Imag'}: fine first layer, coarse second layer.
  \item Line 226: \code{Scattering\_Fourth\_Order\_Imag\_1d(J, 1, filters, filters[:,:-1])}
        — term \code{'Scattering\_Fourth\_Order\_Imag\_Q1'}: both layers coarse.
\end{itemize}
\verified

\paragraph{Normalization.} \code{self.norm}$[s] = \sigma[s]$ (exponent $0.5$,
line 508). \verified — already matched \eqref{eq:paper-norm} before this
session; not modified.

\bigskip
\hrule
\bigskip

\section{\code{Scattering\_Fourth\_Order\_Mod2\_Real\_1d}}
\label{sec:s4mod2real}
\textbf{Code:} \code{potentials\_1d.py:522--616} (post-fix). \quad
\textbf{Term:} \code{'Scattering\_Fourth\_Order\_Mod2\_Real\_Q1'} (only the
$Q=1$ wiring exists — there is no plain, fine-$Q$ \code{'Scattering\_Fourth\_Order\_Mod2\_Real'}
term).

\paragraph{Formula.} Identical to Section~\ref{sec:s4real}, except the envelope
is the \textbf{modulus squared} instead of the modulus:
\begin{equation}
  E^{(2)}[s](t) := |\Wx^{\text{norm}}[s](t)|^2, \qquad
  S_4^{\text{Mod2,Re}}[s_1,s_2,j] = \E_t\Big[\,\mathrm{Re}\big(W_j E^{(2)}[s_1](t)\cdot \overline{W_j E^{(2)}[s_2](t)}\big)\Big].
\end{equation}
$j$ includes the low-pass channel (same default as the non-Mod2 Real class).

\paragraph{Filter bank.} \code{potentials\_builder.py:222}: \code{Scattering\_Fourth\_Order\_Mod2\_Real\_1d(J, 1, filters, filters[:,:-1])}
— \textbf{coarse} ($Q=1$) bank for both layers, explicitly. \verified

\paragraph{Normalization.} \code{self.norm}$[s]=\sigma[s]$ \fixed — was
\code{self.norm}$[s]=\big(\E[|\Wx[s]|^2]\big)^{-1/4}$ before this session (old
line kept commented out). Squaring the envelope changes how many powers of the
normalization factor propagate into the final statistic (two, instead of one),
but this cancels exactly against the extra power picked up from squaring
$|\Wx^{\text{norm}}|$ itself — repeating the scale-tracking argument from
Section~\ref{sec:s4real} for the squared envelope gives a raw-statistic scaling
of $E_{s_1}^{1-2p}E_{s_2}^{1-2p}$, which \emph{also} vanishes only at $p=1/2$.
So the same exponent, $0.5$, is correct for both the modulus and modulus-squared
families — the previous \code{-1/4} value was not merely "right for one, wrong
for the other," it did not achieve scale-invariance for \emph{either}.

\bigskip
\hrule
\bigskip

\section{\code{Scattering\_Fourth\_Order\_Mod2\_Imag\_1d}}
\textbf{Code:} \code{potentials\_1d.py:618--729}. \quad
\textbf{Term:} \code{'Scattering\_Fourth\_Order\_Mod2\_Imag\_Q1'} (again, only
the $Q=1$ wiring exists).

\paragraph{Formula.} As Section~\ref{sec:s4mod2real}, imaginary part, $s_1<s_2$
strictly, $j$ excludes the low-pass channel (same \code{include\_lowpass=False}
convention as the non-Mod2 Imag class):
\begin{equation}
  S_4^{\text{Mod2,Im}}[s_1,s_2,j] = \E_t\Big[\,\mathrm{Im}\big(W_j E^{(2)}[s_1](t)\cdot \overline{W_j E^{(2)}[s_2](t)}\big)\Big].
\end{equation}

\paragraph{Filter bank.} \code{potentials\_builder.py:227}: \code{Scattering\_Fourth\_Order\_Mod2\_Imag\_1d(J, 1, filters, filters[:,:-1])}
— coarse ($Q=1$) bank, both layers. \verified

\paragraph{Normalization.} \code{self.norm}$[s]=\sigma[s]$ (exponent $0.5$,
comment in code: ``matches \code{Imag\_1d} convention (not the -1/4 used by the
Real classes)''). \verified — already correct before this session; not modified.

\bigskip
\hrule
\bigskip

\section{\code{Scalar}}
\textbf{Code:} \code{potentials\_1d.py:731--877}. \quad
\textbf{Terms:} \code{'Scalar\_phi\_quantile\_confine'}, \code{'Scalar\_morlet\_quantile\_confine'},
\code{'Scalar\_psi\_quantile\_confine'}.

\paragraph{Formula.} Not a moment — a \emph{soft empirical CDF} of $|\Wx[\lambda]|$
at $K$ (\code{num\_centers}, default 20) quantile-spaced locations per channel
$\lambda$, i.e. it constrains the shape of the marginal distribution rather than
a single statistic. Calibrated once in \code{fit(x)}: centers
$c_k[\lambda] = Q_{q_k}\big(|\Wx[\lambda](t)|\big)$ placed at $K$ quantiles
$q_k \in [\code{quantile\_min},\code{quantile\_max}]$ (default $[0,1]$), and
widths $s_k[\lambda]$ set from consecutive-quantile spacing times
\code{stride\_sigmas} (default $0.75$). Then, per center $k<K$:
\begin{equation}
  \Sigma[k,\lambda] = \frac{1}{c_k[\lambda]}\;\E_t\left[\,\sigma\!\left(\frac{|\Wx[\lambda](t)| - c_k[\lambda]}{s_k[\lambda]}\right)\right],
\end{equation}
and for $k=K$ (the top quantile, when \code{confine=True}), an extra
polynomial tail correction is added: letting $u(t)=(|\Wx[\lambda](t)|-c_K[\lambda])/s_K[\lambda]$,
the weight becomes $\sigma(u(t))\cdot\big(1+u(t)^8\big)$ before the same
time-average and $1/c_K[\lambda]$ normalization. This is a genuinely different
statistic family from every other class in this file — a histogram/CDF
constraint, not a moment.

\paragraph{Filter bank.} Three wirings, one per bank:
\begin{itemize}
  \item Line 142: \code{Scalar(filters\_Phi, ...)} — term \code{'Scalar\_phi\_quantile\_confine'}: the separate $J$-many low-pass stack.
  \item Line 145: \code{Scalar(filters, ...)} — term \code{'Scalar\_morlet\_quantile\_confine'}: coarse ($Q=1$) bank.
  \item Line 148: \code{Scalar(filters\_Q, ...)} — term \code{'Scalar\_psi\_quantile\_confine'}: fine bank.
\end{itemize}
\verified

\paragraph{Normalization.} No \code{fit\_micro}/\code{self.norm} mechanism at
all — normalization here is implicit in the construction itself: the centers
and widths $(c_k[\lambda], s_k[\lambda])$ are quantiles of the \emph{same}
channel's own data, so the statistic is comparable across scales by
construction, and the explicit $1/c_k[\lambda]$ division further rescales each
bin by its own center. Not directly comparable to \eqref{eq:paper-norm}'s
convention (there is no single "physical units" this class shares with the
moment-based potentials above), but not obviously deficient for its own purpose
either — a different question from the one this note otherwise focuses on.

\bigskip
\hrule
\bigskip

\section{\code{L2p\_norm}}
\label{sec:l2p}
\textbf{Code:} \code{potentials\_1d.py:882--917}. \quad
\textbf{Terms:} \code{'L\_2'} ($p{=}1$), \code{'L\_4'} ($p{=}2$), \code{'L\_6'},
\code{'L\_6\_psi'} ($p{=}3$), \code{'L\_8'} ($p{=}4$; also \code{'L\_9'} and
\code{'L\_10'} — a pre-existing bug overwrites key \code{'L\_8'} instead of
adding a new one, so at most one of \code{L\_8}/\code{L\_9}/\code{L\_10} ever
survives in the returned dict), plus \code{'L\_2\_phi'}...\code{'L\_8\_phi'} and
\code{'L\_2\_lowpass'}, \code{'L\_4\_lowpass'}.

\paragraph{Formula.} A generic, arbitrary-order \emph{single-scale} even moment
of the modulus:
\begin{equation}
  L_{2p}[\lambda] = \frac{\E_t\big[\,|\Wx[\lambda](t)|^{2p}\,\big]}{\code{self.norm}[\lambda]}.
\end{equation}
No cross-scale term of any kind — see the comparison with \code{Scattering\_Third\_Order\_Real\_1d}
in the box below.

\paragraph{Filter bank.} All plain \code{L\_2}...\code{L\_10} use \textbf{coarse}
\code{filters} ($Q=1$); the one exception is \code{'L\_6\_psi'}, which uses
\textbf{fine} \code{filters\_Q} instead (line 106) — same $p=3$ formula, finer
scale grid. The \code{\_phi} variants use \code{filters\_Phi} (the separate
$J$-many low-pass stack); \code{L\_2\_lowpass}/\code{L\_4\_lowpass} slice
\code{filters} down to just its bundled low-pass channel,
\code{filters[:,-1:,:]}. \verified — in your currently-active experiment
configs, only \code{'L\_2\_lowpass'} is actually used (\code{gaussian\_theta\_scaling.ipynb}); none of the others appear in your current \code{terms} lists.

\paragraph{Normalization.} \code{self.norm}$[\lambda] = \big(\E[|\Wx[\lambda]|^{2p}]\big)^{-1/2}$
(fit from reference data), giving
$L_{2p}[\lambda] = \E_t[|\Wx[\lambda]|^{2p}] \cdot M_{\mathrm{ref}}[\lambda]^{1/2}$
— on matching data this is $\approx M_\mathrm{ref}[\lambda]^{3/2}$, \emph{not}
$1$. \suspect\ Two independent problems, not one:
\begin{enumerate}
  \item Even granting the class's own apparent goal (self-normalize each order
        to $\approx 1$ on its reference data), the active exponent doesn't
        achieve it — it \emph{inflates} rather than divides down, the same
        class of mistake found and fixed in Sections~\ref{sec:s4real}--\ref{sec:s4mod2real}.
        A commented-out alternative sits right below the active line
        (\code{self.norm = raw\_moment}, no exponent), which \emph{would}
        self-normalize correctly.
  \item Even that corrected version would not implement \eqref{eq:paper-norm}:
        it normalizes $L_{2p}$ by \emph{its own} $2p$-th moment
        $M_{2p}[\lambda]=\E[|\Wx[\lambda]|^{2p}]$, not by $\sigma[\lambda]^{2p}=\big(\E[|\Wx[\lambda]|^2]\big)^p$
        as the paper's convention requires. These coincide only at $p=1$ (that
        is exactly \code{L\_2}); for $p=2,3,4$ they are genuinely different
        quantities (e.g. for a Gaussian process, $\E[|\Wx|^4]=2(\E[|\Wx|^2])^2$,
        so \code{L\_4}'s reference differs from $\sigma^4$ by more than a
        relabeling). Self-normalizing each order against its own value also
        forfeits the paper's actual point: keeping every renormalized order in
        the \emph{same} units ($\sigma$) so their magnitudes stay comparable
        to each other, not just individually $\approx 1$.
\end{enumerate}
Not fixed in this session — flagged for a separate decision.

\bigskip
\hrule
\bigskip

\section{\code{L2p1\_norm}}
\label{sec:l2p1}
\textbf{Code:} \code{potentials\_1d.py:919--954}. \quad
\textbf{Terms:} \code{'L\_1'} ($p{=}0$), \code{'L\_3'} ($p{=}1$), \code{'L\_5'}
($p{=}2$), \code{'L\_7'} ($p{=}3$), plus \code{'L\_1\_phi'}...\code{'L\_7\_phi'}.

\paragraph{Formula.} The odd-order counterpart of Section~\ref{sec:l2p}:
\begin{equation}
  L_{2p+1}[\lambda] = \frac{\E_t\big[\,|\Wx[\lambda](t)|^{2p+1}\,\big]}{\code{self.norm}[\lambda]}.
\end{equation}
Again purely single-scale; no term in this class can represent a cross-scale
(order-3-or-higher scattering) coefficient — see the boxed note below.

\paragraph{Filter bank.} All plain \code{L\_1}...\code{L\_7} use \textbf{coarse}
\code{filters} ($Q=1$); the \code{\_phi} variants use \code{filters\_Phi}. There
is no fine-bank ($Q{>}1$) wiring for this class anywhere in
\code{get\_1d\_potentials}. \verified — none of these terms appear in your
currently-active experiment configs.

\paragraph{Normalization.} \code{self.norm}$[\lambda] = \E[|\Wx[\lambda]|^{2p+1}]$
(the raw reference moment itself, no exponent at all). This \emph{does}
self-normalize $L_{2p+1}$ to $\approx 1$ on matching data — but that is dividing
a statistic by \emph{itself}, which is a different (weaker) claim than the
paper's $/\sigma^{2p+1}$ convention, for the same reason given in point 2 of
Section~\ref{sec:l2p}: it does not keep different orders in comparable, shared
units. \suspect

\bigskip
\hrule
\bigskip

\section*{Summary table}
\begin{center}
\renewcommand{\arraystretch}{1.3}
\begin{tabular}{@{}p{4.6cm}p{2.4cm}p{2.6cm}p{4.7cm}@{}}
\toprule
\textbf{Class} & \textbf{Order} & \textbf{Filter bank(s)} & \textbf{Normalization} \\
\midrule
\code{Scattering\_First\_Order\_1d} & 1 & fine, band-pass only & none \\
\code{Scattering\_Second\_Order\_1d} & 2 & fine \emph{or} coarse & none \\
\code{Scattering\_Second\_Order\_Bulk\_1d} & 2 (trimmed) & fine & tail-window calib., not \eqref{eq:paper-norm} \\
\code{Scattering\_Third\_Order\_Real\_1d} & 3 (cross-scale) & fine, \textbf{but only correct at $Q{=}1$} & none \\
\code{Scattering\_Fourth\_Order\_Real\_1d} & 4 (cross-scale) & fine+coarse, or coarse only & $\sigma[s]^{0.5}$ -- fixed \\
\code{Scattering\_Fourth\_Order\_Imag\_1d} & 4 (cross-scale) & fine+coarse, or coarse only & $\sigma[s]^{0.5}$ -- already correct \\
\code{Scattering\_Fourth\_Order\_Mod2\_Real\_1d} & 4, mod-squared & coarse only ($Q{=}1$) & $\sigma[s]^{0.5}$ -- fixed \\
\code{Scattering\_Fourth\_Order\_Mod2\_Imag\_1d} & 4, mod-squared & coarse only ($Q{=}1$) & $\sigma[s]^{0.5}$ -- already correct \\
\code{Scalar} & CDF/quantiles & phi-stack, coarse, or fine & implicit in quantile calib. \\
\code{L2p\_norm} & even, $2,4,6,8$ & coarse (mostly), one fine & wrong exponent \emph{and} wrong reference quantity \\
\code{L2p1\_norm} & odd, $1,3,5,7$ & coarse only & self-normalizes, not $\sigma$-based \\
\bottomrule
\end{tabular}
\end{center}

\end{document}
